\documentclass[../../../main]{subfiles}

\ifSubfilesClassLoaded{%
    \zexternaldocument*{../1.Category-Theory/exported/aux/main}%
    \zexternaldocument*{../1.Additive-and-abelian-categories/exported/aux/main}%
    \zexternaldocument*{../2.Complexes-and-cohomology/exported/aux/main}%
    \zexternaldocument*{../3.Cones-and-homotopy/exported/aux/main}%
    \zexternaldocument*{1.Injective-and-projective-objects}%
}{}

\begin{document}
\section{Resolutions and homotopy}

\begin{definition}~ \label{def:injective-projective-resolution}
  \begin{itemize}
    \item An \textbf{injective resolution} of an object \(A \in \category{A}\) is a \hyperref[def:quasi-isomorphism]{quasi-isomorphism} \(\iota(A) \to Q^{\bullet}\), where \(Q^{\bullet}\) is a \hyperref[def:complex]{complex} in \(\Complex[\geq 0]{I}\).
    \item A \textbf{projective resolution} of an object \(A \in \category{A}\) is a \hyperref[def:quasi-isomorphism]{quasi-isomorphism} \(P^{\bullet} \to \iota(A)\), where \(P^{\bullet}\) is a \hyperref[def:complex]{complex} in \(\Complex[\leq 0]{P}\).
  \end{itemize}

  We will refer to this resolutions using \(Q^{\bullet}, P^{\bullet}\), \textit{forgetting} about the \hyperref[def:quasi-isomorphism]{quasi-isomorphism}.
\end{definition}

\begin{lemma}
  Let \(\category{A}\) be an \hyperref[def:abelian-category]{abelian category}.
  \begin{itemize}
    \item If \(\category{A}\) has \hyperref[def:enough-injectives-projectives]{enough injectives} then every object \(A \in \category{A}\) admits an \hyperref[def:injective-projective-resolution]{injective resolution}.
    \item If \(\category{A}\) has \hyperref[def:enough-injectives-projectives]{enough projectives} then every object \(A \in \category{A}\) admits a \hyperref[def:injective-projective-resolution]{projective resolution}.
  \end{itemize}
\end{lemma}
\begin{proof}
  We will do the \hyperref[def:enough-injectives-projectives]{enough injectives} case.

  Since \(\category{A}\) has \hyperref[def:enough-injectives-projectives]{enough injectives} there exists an \hyperref[def:injective-projective-object]{injective object} \(Q^{0} \in \category{A}\) and a \hyperref[def:mono-epi-morphism]{monomorphism} \(d^{0}: A \monoarrow Q^{0}\).
  Then there exists an \hyperref[def:injective-projective-object]{injective object} \(Q^{1} \in \category{A}\) an a \hyperref[def:mono-epi-morphism]{monomorphism} \(d^{1}: \coker  d^{0} \monoarrow Q^{0}\) and so on.
  Pictorially
  \begin{center}
    \begin{tikzcd}
      A \arrow[r, tail, dotted, "d^{0}"] & Q^{0} \arrow[d, two heads] \\
      & \coker d^{0} \arrow[r, tail, dotted] & Q^{1} \arrow[d, two heads] \\
      & & \coker d^{1} \arrow[r, tail, dotted] & Q^{2} \arrow[d, two heads] \\
      & & & \vdots
    \end{tikzcd}
  \end{center}
  Therefore following this diagram we get an \hyperref[def:injective-projective-resolution]{injective resolution}
  \begin{center}
    \begin{tikzcd}
      Q^{\bullet} = \cdots \arrow[r] & 0 \arrow[r] & Q^{0} \arrow[r] & Q^{1} \arrow[r] & Q^{2} \arrow[r] & \cdots
    \end{tikzcd}
  \end{center}

\end{proof}

We will assume that \(\category{A}\) is an \hyperref[def:abelian-category]{abelian-cateogory} with \hyperref[def:enough-injectives-projectives]{enough injectives}.

By \textit{flipping} the corresponding arrows, all the results may be adapted to the \hyperref[def:injective-projective-object]{projective} case.

\begin{lemma}\label{lemma:lifting-to-cohomology}
  Let \(A\) be an object in a \(\category{A}\), let \(M^{\bullet}\) be an \hyperref[def:injective-projective-resolution]{injective resolution} of \(A\) and let \(Q^{\bullet}\) be a \hyperref[def:complex]{complex} in \(\Complex[\geq 0]{I}\).
  Let \(\varphi: \cohomology{0}(M^{\bullet}) = A \to \cohomology{0}(Q^{\bullet})\) be an arbitrary morphism. Then
  \begin{itemize}
    \item There exists a \hyperref[def:complex-morphism]{morphism of complexes} \(\alpha: M^{\bullet} \to Q^{\bullet}\) inducing \(\varphi\) at the level of \(\cohomology{0}\).
    \item Different morphisms satisfying the above are \hyperref[def:homotopy-equivalence]{homotopy equivalent}.
  \end{itemize}
\end{lemma}
\begin{proof}
  We have to define \(\forall i \in \mathbb{Z} \quad \alpha^{i}: M^{i} \to Q^{i}\).
  Since for \(i < 0\) we have \(\alpha^{i} = 0\) we can truncate the complexes, obtaining
  \begin{center}
    \begin{tikzcd}
      0 \arrow[r] & \cohomology{0}(M^{\bullet}) = A \arrow[r, "\varepsilon"] \arrow[d, "\varphi"] & M^{0} \arrow[r, "d_{M^{\bullet}}^{0}"] \arrow[d, dotted, "\alpha^{1}"] & M^{1} \arrow[r, "d_{M^{\bullet}}^{1}"] \arrow[d, dotted, "\alpha^{1}"] & M^{2} \arrow[r, "d_{M^{\bullet}}^{2}"] \arrow[d, dotted, "\alpha^{2}"] & \cdots \\
      0 \arrow[r] & \cohomology{0}(Q^{\bullet}) \arrow[r, "\iota"'] & Q^{0} \arrow[r, "d_{Q^{\bullet}}^{0}"'] & Q^{1} \arrow[r, "d_{Q^{\bullet}}^{1}"'] & Q^{2} \arrow[r, "d_{Q^{\bullet}}^{2}"'] & \cdots
    \end{tikzcd}
  \end{center}
  Note that the top \hyperref[def:complex]{complex} is \hyperref[def:exact-complex]{exact}.

  We will construct \(\alpha^{\bullet}\) using induction:
  \begin{itemize}
    \item Case \(n = 0\). Consider the following diagram
          \begin{center}
            \begin{tikzcd}
              A \arrow[r, tail, "\varepsilon"] \arrow[d, "\varphi"] & M^{0} \\
              \cohomology{0}(Q^{\bullet}) \arrow[r, "\iota"'] & Q^{0}
            \end{tikzcd}
          \end{center}
          Since \(Q^{0}\) is \hyperref[def:injective-projective-object]{injective} using its universal property with the morphism \(\iota \circ \varphi\) we obtain morphism \(\alpha^{0}: M^{0} \to Q^{0}\) making the diagram commute.
    \item Inductive step. Suppose we have already construced \(\alpha^{0}, \ldots, \alpha^{n-1}\) such that
          \begin{align*}
            \forall k \in \{1, \ldots, n-1\} \quad \alpha^{k} \circ d_{M^{\bullet}}^{k-1} = d_{Q^{\bullet}}^{k} \circ \alpha^{k-1}
          \end{align*}
          We want to construct \(\alpha^{n}: M^{n} \to Q^{n}\).

          Consider the diagram
          \begin{center}
            \begin{tikzcd}
              M^{n-1} \arrow[r, "d_{M^{\bullet}}^{n-1}"] \arrow[d, "\alpha^{n-1}"']
              & M^{n}  \\
              Q^{n-1} \arrow[r, "d_{Q^{\bullet}}^{n-1}"']
              & Q^{n}
            \end{tikzcd}
          \end{center}
          Let \(f = d_{Q^{\bullet}}^{n-1} \circ \alpha^{n-1}: M^{n-1} \to Q^{n}\).
          Note that
          \begin{align*}
            f \circ d_{M^{\bullet}}^{n-2} = d_{Q^{\bullet}}^{n-1} \circ \alpha^{n-1} \circ d_{M^{\bullet}}^{n-2} =  \underbrace{d_{Q^{\bullet}}^{n-1} \circ d_{Q^{\bullet}}^{n-2}}_{0} \circ \alpha^{n-2} = 0
          \end{align*}
          meaning that \(f\) vanishes on \(\im d_{M^{\bullet}}^{n-2} = \ker d_{M^{\bullet}}^{n-1}\) by \hyperref[def:exact-complex]{exactness} of \(M^{\bullet}\).

          Therefore \(f\) induces a map \(\overline{f}: \im d_{M^{\bullet}}^{n-1} \to Q^{n}\)
          \begin{center}
            \begin{tikzcd}
              M^{n-1} \arrow[r] \arrow[rr, bend left, "d_{M^{\bullet}}^{n-1}"] \arrow[d, "\alpha^{n-1}"'] \arrow[drr, "f"'] & \im d_{M^{\bullet}}^{n-1} \arrow[r, tail] \arrow[dr, dashed, "\overline{f}"] & M^{n} \\
              Q^{n-1} \arrow[rr, "d_{Q^{\bullet}}^{n-1}"'] & & Q^{n}
            \end{tikzcd}
          \end{center}
          Using that \(Q^{n}\) is \hyperref[def:injective-projective-object]{injective} we obtain a morphism \(\alpha^{n}: M^{n} \to Q^{n}\).
  \end{itemize}
  \todo[inline]{Falta unicitat}
\end{proof}

\begin{proposition}\label{prop:injective-resolutions-are-homotopy-equivalent}
  Any two \hyperref[def:injective-projective-resolution]{injective resolutions} (resp. \hyperref[def:injective-projective-resolution]{projective}) of an object \(A\) in an \hyperref[def:abelian-category]{abelian category} \(\category{A}\) are \hyperref[def:homotopy-equivalence]{homotopy equivalent}.
\end{proposition}
\begin{proof}
  Let \(I^{\bullet}, J^{\bullet}\) be two \hyperref[def:injective-projective-resolution]{injective resolutions} of an object \(A \in \category{A}\).
  By definition we have \(\cohomology{0}(I^{\bullet}) = A = \cohomology{0}(J^{\bullet})\).

  Consider the identity morphism \(\id_{A}: \cohomology{0}(I^{\bullet}) \to \cohomology{0}(J^{\bullet})\).
  By \autoref{lemma:lifting-to-cohomology} taking \(M^{\bullet} = I^{\bullet}\) and \(Q^{\bullet} = J^{\bullet}\) there exists a \hyperref[def:complex-morphism]{morphism of complexes} \(f: I^{\bullet} \to J^{\bullet}\) that induces \(\id_{A}\) at the level of \(\cohomology{0}\).
  Taking now \(M^{\bullet} = J^{\bullet}\) and \(Q^{\bullet} = I^{\bullet}\) we find another \hyperref[def:complex-morphism]{morphism of complexes} \(g: J^{\bullet} \to I^{\bullet}\) inducing \(\id_{A}\) at the level of \(\cohomology{0}\).

  If we consider the composite morphism \(g \circ f: I^{\bullet} \to I^{\bullet}\) we have that
  \begin{equation*}
    \cohomology{0}(g \circ f) = \cohomology{0}(g) \circ \cohomology{0}(f) = \id_{A} \circ \id_{A} = \id_{A}
  \end{equation*}
  Therefore since \(\id_{I^{\bullet}}: I^{\bullet} \to I^{\bullet}\) also induces \(\id_{A}\) at the level of \(\cohomology{0}\) by the second part of \autoref{lemma:lifting-to-cohomology}
  \begin{equation*}
    f \circ g \simeq \id_{I^{\bullet}}
  \end{equation*}
  If we consider \(f \circ g\) analogously we get
  \begin{equation*}
    g \circ f \simeq \id_{J^{\bullet}}
  \end{equation*}
  Therefore \(f\) and \(g\) are an \hyperref[def:homotopy-equivalence]{homotopy equivalence}.
\end{proof}


\begin{corollary} \label{cor:injective-have-same-cohomology}
  Any two \hyperref[def:injective-projective-resolution]{injective resolutions} (resp. \hyperref[def:injective-projective-resolution]{projective}) of an object \(A\) in an \hyperref[def:abelian-category]{abelian category} \(\category{A}\) have \hyperref[def:categorical-isomorphism]{isomorphic} \hyperref[def:cohomology]{cohomology}.
\end{corollary}
\begin{proof}
  Direct from \autoref{prop:injective-resolutions-are-homotopy-equivalent} and \autoref{cor:homotopy-equivalent-isomorphic-cohomology}
\end{proof}


\begin{theorem}\label{thm:existence-of-lifting}
  Let \(A_{0}, A_{1}\) be two objects in an \hyperref[def:abelian-category]{abelian category} \(\category{A}\).
  Let \(I_{i}^{\bullet}\) be an \hyperref[def:injective-projective-resolution]{injective resolution} of \(A_{i}, i \in \{0, 1\}\).
  Then every morphism \(\varphi: A_{0} \to A_{1}\) in \(\category{A}\) is induced by a morphism \(\alpha^{\bullet}: I_{0}^{\bullet} \to I_{1}^{\bullet}\), uniquely determined up to \hyperref[def:homotopy]{homotopy}.
  We call the morphism \(\alpha^{\bullet}\) a \textit{lifting} of \(\varphi\).
\end{theorem}
\begin{proof}
  We will first construct \(\alpha^{\bullet}\) by induction.
  
  We have \hyperref[def:injective-projective-resolution]{injective resolutions}
  \begin{center}
    \begin{tikzcd}
      I_{0}^{\bullet} = 0 \arrow[r] & A_{0} \arrow[r, "\varepsilon_{0}"] & I_{0}^{0} \arrow[r, "d_{0}^{0}"] & \cdots \\
      I_{1}^{\bullet} = 0 \arrow[r] & A_{1} \arrow[r, "\varepsilon_{1}"] & I_{1}^{0} \arrow[r, "d_{1}^{0}"] & \cdots 
    \end{tikzcd}
  \end{center}
  Consider \(\varepsilon_{1} \circ \varphi: A_{0} \to I_{1}^{0}\).
  Since \(\varepsilon_{0}: A_{0} \to I_{0}^{0}\) is a \hyperref[def:mono-epi-morphism]{monomorphism} and \(I_{1}^{0}\) is \hyperref[def:injective-projective-object]{injective} there exists a morphism \(\alpha_{0}: I_{0}^{0} \to I_{1}^{1}\) extending \(\varepsilon_{1} \circ \varphi\)
  \begin{center}
    \begin{tikzcd}
      A_{0} \arrow[r, "\varepsilon_{1} \circ \varphi"] \arrow[d, "\varepsilon_{0}"', tail] & I_{1}^{0} \\
      I_{0}^{0} \arrow[ur, "\widetilde{\alpha}_{0}"', dotted]
    \end{tikzcd}
  \end{center}
  Therefore we can define \(\alpha_{0} \coloneq \widetilde{\alpha}_{0} \circ \varepsilon_{0}\).
  
  Supose now that for all \(i \in \{0, \ldots, n\}\) we have a morphism \(\alpha_{i}: I_{0}^{i} \to I_{1}^{i}\) such that \(d_{1}^{i-1} \circ \alpha^{i-1} = \alpha_{i} \circ d_{0}^{i}\). We want to construct \(\alpha_{n}: I_{0}^{n+1} \to I_{1}^{n+1}\).

  Consider the morphism \(d_{1}^{n} \circ \alpha^{n}: I_{0}^{n} \to I_{1}^{n+1}\).
  We have that
  \begin{equation*}
    (d_{1} \circ \alpha^{n}) \circ d_{0}^{n-1} = d_{1}^{n} \circ (\alpha^{n} \circ d_{0}^{n-1}) = d_{1}^{n} \circ d_{1}^{n-1} \circ \alpha^{n-1} = 0 \circ \alpha^{n-1} = 0
  \end{equation*}
  Therefore \(d_{1}^{n} \circ \alpha^{n}\) factors through \(\coker d_{0}^{n-1} \cong \operatorname{im} d_{0}^{n} \subseteq I_{0}^{n+1}\) and since \(I_{1}^{n+1}\) is \hyperref[def:injective-projective-object]{injective}
  \begin{center}
    \begin{tikzcd}
        I_{0}^{n+1} \arrow[r, "d_{1}^{n} \circ \alpha^{n}"] \arrow[d, "d_{0}"', tail] & I_{1}^{n+1} \\
       \operatorname{im} d_{0}^{n} \arrow[ur, "\widetilde{\alpha}_{n}"', dotted]
    \end{tikzcd}
  \end{center}
  Defining \(\alpha_{n} \coloneq d_{0} \circ \widetilde{\alpha}_{n}\) we finish the proof of existence.
  
  To check uniqueness, consider \(\beta^{\bullet}: I_{0}^{\bullet} \to I_{1}^{\bullet}\) another lifting of \(\varphi\).
  Let \(\gamma^{\bullet} \coloneq \alpha^{\bullet} - \beta^{\bullet}\), which is a lifting of \(0: A_{0} \to A_{1}\).
  We are going to construct \hyperref[def:homotopy]{homotopy} \(h^{n}: I_{0}^{n} \to I_{1}^{n-1}\) such that
  \begin{equation*}
    \gamma^{n} = d_{1}^{n-1} \circ h^{n} + h^{n+1} \circ d_{0}^{n}
  \end{equation*}
  by induction, where \(h^{0} \coloneq 0\) and \(I_{1}^{-1} = A_{1}\).

  For the base case, since \(\gamma^{0} \circ \varepsilon_{0} = 0\) we have that \(\gamma^{0}\) factors through \(\coker \varepsilon_{0} \cong \operatorname{im} d_{0}^{0} \subseteq I_{0}^{1}\).
  Therefore since \(I_{0}^{1}\) is \hyperref[def:injective-projective-object]{injective} we have
  \begin{center}
    \begin{tikzcd}
      I_{0}^{0} \arrow[r, "\gamma^{0} \circ \varepsilon_{0}"] \arrow[d, "d_{0}"', tail] & I_{0}^{1} \\
      \operatorname{im} d_{0}^{0} \arrow[ur, "\widetilde{h}_{1}"', dotted]
    \end{tikzcd}
  \end{center}
  \todo[inline]{FINISH}
\end{proof}
\end{document}
